\scanpage{77}
\mspage{38}
dans ce cas le critère homologique s'avère insuffisant\ldots) Il faudrait donc calculer cet \deleted{élément de}\inserted{homomorphisme}
\[
(11)\qquad T_\ell\bigl(\simeq H_1(\eta,\mathbb Z_\ell)\bigr)\longrightarrow T_\ell^I/T_\ell
\]
pour constater qu'il n'est pas nul dans le cas envisagé, $\nu\geq2$ (et traiter séparément le cas $\nu=1$).

\deleted{[On va faire des calculs intrinsèques, en
\[
\unsure{\operatorname{Ker}}(\mathcal J^{\unsure{\bullet|0}}_{U/S}\longrightarrow\mathcal J^{\unsure{\bullet|0}}_{X/S})\simeq G_m^I/G_m
\]
\[
H_1(U,\mathbb Z_\ell)\simeq H_1(\mathcal J^1_{U/S},\mathbb Z_\ell)\underset{\text{isom de Kummer}}{\simeq}\lim_{\longleftarrow}\deleted{\unsure{\ell^n}}{}_{\ell^n}(\mathcal J^0_{U/S})
\]
On considère donc l'image inverse par $\eta\longrightarrow\mathcal J^1_{U/S}$ (composé de $\eta\longrightarrow U\longrightarrow\mathcal J^1_{U/S}$) des isogénies}

Je vais deviner le résultat :

Soit \deleted{\unsure{q}} $x=g_{i_0}(s)$, $A=\underline O_{X,x}$, $V$ l'anneau de $S$, $\mathcal J_{i_0}$ l'idéal de l'hom $A\xrightarrow{g_{i_0}^*}V$ associé à $g_{i_0}$, c'est donc un idéal inversible de $A$ -- soit de $\unsure{\overline m|\overset{m}{=}}$ $\mathcal J_f$ l'idéal associé à $f^*:A\longrightarrow V$, et considérons $g_{i_0}^*(\mathcal J_f)$, c'est un idéal de $V$ engendré par un générateur, et comme $g_{i_0}\neq f$, on voit que cet idéal n'est pas nul. Soit $\mu=\operatorname{long}V/g_{i_0}^*(\mathcal J_f)\ \deleted{\unsure{=}}$, cet entier
\margin{76}
